Reinforcement learning against a verifier improves reasoning models. Whether it creates capabilities the base lacked
or raises the probability of behaviour the base already had is contested, and large-model studies can rarely change
what the base saw. We study the question in a 9.56M-parameter transformer prover, trained from scratch on 155,000
synthetic natural-deduction proofs of propositional logic, judged only by Lean~4, with exact likelihoods for any
proof. We remove one inference composition from pretraining, the classical reductio \sig{} (double-negation
elimination applied to a refuted negation), refill the data to the same size, and run 16 rounds of expert iteration.
In a pre-registered test, no \sig-free base acquired a classical theorem family (0/6 seeds), against 6/6
full-exposure seeds (one-sided Boschloo $p=0.00024$). Counting only families acquired during RL: 5/6 vs 0/6
($p=0.0032$). Restoring \sig{} in 1\% of pretraining restored acquisition on 3/3 seeds. The knockout bases were not
\sig-free: four of six wrote \sig{} proofs of held-out family members (1--6 in 1.1--1.4 million attempts), only two
are certified out of reach at RL's compute budget, and the one knockout run that made progress did so mostly on
members its base could already reach. Measurement matters too: of the 51--60 theorems per seed that RL
``creates'' at equal sample counts, 14--21 remain once the base gets RL's compute as attempts, under every budget
allocation tried. In this organism, expert iteration amplified a composition pretraining had exposed, even at 1\%,
but did not build it within 16 rounds from a base that had never seen it.
